\section{Why $2ab$ measures both height and lifetime}
\label{sec:lifetime-law}
%======================================================================

The same imaginary coordinate that measures factor imbalance also measures how long a source remains in transport.  This is the key space--time link in the construction.

Define the signed continuous delay
\begin{equation}
 \tau(c,q):=q-\sqrt{cq}.
 \label{eq:tau}
\end{equation}
It is positive exactly for transport-capable points, zero at $c=q$, and negative for points that enter already smooth.

Put
\begin{equation}
 r:=\frac cq.
 \label{eq:r-ratio}
\end{equation}
Then
\begin{equation}
 Y=\frac{q^2}{2}(1-r^2)
 \label{eq:Y-r}
\end{equation}
and
\begin{equation}
 \tau=q(1-\sqrt r).
 \label{eq:tau-r}
\end{equation}
Since
\begin{equation}
 1-r^2=(1-\sqrt r)(1+\sqrt r)(1+r),
 \label{eq:factor-r}
\end{equation}
we obtain an exact identity.

\begin{theorem}[Exact height--lifetime identity]
\label{thm:exact-height-lifetime}
For every positive pair $(c,q)$,
\begin{equation}
 \boxed{
 Y
 =
 \frac{q\tau}{2}
 \left(1+\sqrt{\frac cq}\right)
 \left(1+\frac cq\right).
 }
 \label{eq:Y-tau-exact}
\end{equation}
Equivalently,
\begin{equation}
 \boxed{
 \tau
 =
 \frac{2Y}
 {q\left(1+\sqrt{c/q}\right)\left(1+c/q\right)}.
 }
 \label{eq:tau-Y-exact}
\end{equation}
\end{theorem}

The sign of $Y$ is therefore exactly the sign of the continuous delay.  On the transport side $0<c/q<1$, the denominator in \cref{eq:tau-Y-exact} lies between $q$ and $4q$, giving
\begin{equation}
 \frac{Y}{2q}
 \le\tau
 \le\frac{2Y}{q}.
 \label{eq:tau-bounds}
\end{equation}
Since
\begin{equation}
 L(m)=q-1-\lfloor\sqrt m\rfloor
 \le q-\sqrt m=\tau(c,q)
 \label{eq:L-tau}
\end{equation}
for transport points,
\begin{equation}
 \boxed{L(m)\le\frac{2Y_m}{q_m}.}
 \label{eq:L-Y}
\end{equation}

In squared-space coordinates, $q=\sqrt{|w|+\Impart w}$, so
\begin{equation}
 L(w)
 \le
 \frac{2(\Impart w)_+}
 {\sqrt{|w|+\Impart w}}.
 \label{eq:L-w}
\end{equation}
A point close to the real axis cannot remain in transport for many prefix steps.

%======================================================================
